\documentclass[12pt,a4paper]{article}
\usepackage[utf8]{inputenc}
\usepackage[spanish]{babel}
\usepackage{geometry}
\usepackage{listings}
\usepackage{xcolor}
\usepackage{parskip}
\usepackage{amsmath}
\usepackage{hyperref}

\geometry{margin=2.5cm}

\lstset{
  basicstyle=\ttfamily\small,
  keywordstyle=\color{blue}\bfseries,
  commentstyle=\color{gray},
  breaklines=true, frame=single, language=Java,
  numbers=left, numberstyle=\tiny\color{gray}
}

\title{\textbf{Programación Lineal con Backtracking}\\
  \large Documentación Técnica}
\author{Inteligencia Artificial}
\date{}

\begin{document}
\maketitle
\tableofcontents
\newpage

\section{Descripción del Proyecto}
Dos aplicaciones que usan \textbf{Backtracking} para resolver problemas de programación lineal entera (PLE):
\begin{enumerate}
  \item \textbf{BacktrackingPLE}: solucionador genérico de PL entera.
  \item \textbf{BacktrackingCarpinteriaPLE}: problema específico de optimización de ganancias.
\end{enumerate}

\section{Programación Lineal Entera}

La PLE busca maximizar o minimizar una función objetivo lineal sujeta a restricciones lineales, con la condición de que las variables sean \textbf{enteras no negativas}:

\[
\text{Maximizar: } Z = c_1 x_1 + c_2 x_2 + \ldots + c_n x_n
\]
\[
\text{Sujeto a: } a_{ij} x_j \leq b_i \quad \forall i,\quad x_j \geq 0,\ x_j \in \mathbb{Z}
\]

\section{Algoritmo de Backtracking para PLE}

El backtracking explora el espacio de soluciones enteras de forma sistemática:

\begin{lstlisting}
void backtrack(int[] variables, int indice,
               double[][] A, double[] b, double[] c) {
    if (indice == variables.length) {
        if (cumpleRestricciones(variables, A, b)) {
            double z = calcularObjetivo(variables, c);
            if (z > mejorZ) {
                mejorZ = z;
                mejorSolucion = variables.clone();
            }
        }
        return;
    }

    // Calcular límite superior para la variable actual
    int limSup = calcularLimiteSuperior(indice, A, b);

    for (int v = 0; v <= limSup; v++) {
        variables[indice] = v;
        // Poda: si la solución parcial ya no es factible
        if (esParcialmenteFactible(variables, indice,
                                   A, b)) {
            backtrack(variables, indice+1, A, b, c);
        }
    }
    variables[indice] = 0; // restaurar
}
\end{lstlisting}

\subsection{Poda}
La \textbf{poda} elimina ramas que no pueden conducir a una solución óptima:
\begin{itemize}
  \item Si la asignación parcial ya viola alguna restricción, se descarta esa rama.
  \item Si el límite superior de la función objetivo para la solución parcial es menor al mejor encontrado, se poda.
\end{itemize}

\section{Problema de la Carpintería}

\subsection{Formulación}
Una carpintería fabrica mesas ($x_1$) y sillas ($x_2$):

\[
\text{Maximizar: } Z = 5x_1 + 4x_2 \quad \text{(ganancia en \$)}
\]
\[
\text{Sujeto a:}\quad
6x_1 + 4x_2 \leq 24 \quad \text{(madera disponible)}
\]
\[
x_1 + 2x_2 \leq 6 \quad \text{(horas de trabajo)}
\]
\[
x_1, x_2 \geq 0, \quad x_1, x_2 \in \mathbb{Z}
\]

\subsection{Solución óptima}
El backtracking encuentra $x_1 = 3,\ x_2 = \frac{3}{2}$, pero al requerir enteros: $x_1 = 3,\ x_2 = 1$ con $Z = 19$.

\section{Interfaz Gráfica}

\begin{itemize}
  \item Formularios para ingresar coeficientes de la función objetivo y restricciones.
  \item Tabla de restricciones dinámica (se pueden agregar/quitar filas).
  \item Botón \textbf{Resolver} ejecuta el backtracking y muestra la solución óptima.
  \item Muestra los valores de todas las variables y el valor óptimo de $Z$.
\end{itemize}

\end{document}
